\section{Significance of the Study}
Hospitals are essential facilities whose seismic performance has consequences beyond the direct repair of structural damage. By comparing fixed-base and flexible-base representations of the same hospital prototype, this study provides a consistent basis for examining how foundation flexibility changes structural demand in the Nepalese context.

The comparison links modal response, displacement, drift, shear, foundation response, and member-force demand with nonlinear pushover behaviour and preliminary fragility. Considering these measures together is useful because SSI does not necessarily affect every response quantity in the same direction. The study therefore emphasizes deformation and demand redistribution rather than relying on a single indicator such as base shear.

\section{Scope and Limitations of the Study}
The study is limited to the selected G+5 reinforced-concrete hospital prototype and four analytical support conditions: fixed base, hard soil, medium soil, and soft soil. Response-spectrum analysis provides the principal linear comparison, while SAP2000 pushover capacity and the reported FEMA~356 performance-point coordinates provide the nonlinear static comparison. Foundation settlement is evaluated from ETABS results. The fragility assessment uses nonlinear time-history results from twelve ground motions. Because the available results do not specify the response direction and the Barpark/Barpak-labelled record could not be independently verified, the fragility findings are preliminary and should not be generalized to all hospital buildings or all Nepalese site conditions.

\section{Organization of the Study}
The thesis is organized into five chapters. Chapter 1 defines the research problem, questions, objectives, significance, scope, and limitations. Chapter 2 reviews literature on seismic vulnerability, hospital performance, SSI, performance-based assessment, pushover analysis, and fragility. Chapter 3 describes the analytical methodology and verification of the analysis results. Chapter 4 presents and discusses the linear, nonlinear, foundation, and fragility results. Chapter 5 summarizes the principal conclusions, recommendations, and limitations of the study.
